\documentclass[11pt]{article}
\usepackage[T1]{fontenc}
\usepackage[margin=0.92in]{geometry}
\usepackage{amsmath,amssymb}
\usepackage[hidelinks]{hyperref}
\usepackage{microtype}

\newcommand{\C}{\mathbb C}
\newcommand{\Z}{\mathbb Z}
\newcommand{\HW}{\mathrm{HW}}
\newcommand{\op}{\mathrm{op}}

\title{The composite-dimensional Heisenberg--Weyl\\
Bohnenblust--Hille question}
\author{Literature-already-answered note for arXiv:2301.01438}
\date{11 August 2026}

\begin{document}
\maketitle

\paragraph{Status.}
\textbf{Answered in the direct same-author split/follow-up arXiv:2406.08509,}
with a degree loss.  The answer gives a dimension-free inequality for every
composite local dimension.  It does not establish the prime-case exponent and
does not settle optimal growth of the quantum BH constants.

\section*{Source question}
The conclusion of Klein--Slote--Volberg--Zhang \cite{source}, on source PDF
page~32, says that the eigenspace relations needed for qudit
Bohnenblust--Hille inequalities are mostly complete, but asks what can be said
for the Heisenberg--Weyl (HW) basis when \(K\) is composite.  It then asks the
separate question whether exponential dependence of the quantum BH constants
on degree is necessary.

\section*{Direct later answer}
Slote--Volberg--Zhang \cite{answer} is the focused qudit-BH paper by three of
the four source authors.  Its arXiv record explains that an old version had
appeared as arXiv:2301.01438v2 before that record was replaced, and that the
focused version extends the results to more general qudit systems.  This makes
the source relation direct rather than inferential.

Fix a non-prime integer \(K\ge4\), \(d\ge1\), and an observable
\(A\in M_K(\C)^{\otimes n}\) of HW degree at most \(d\), using the degree
convention of \cite{answer}.  Write \(\widehat A\) for its HW coefficient
array.  Theorem~7 of \cite{answer}, on supporting PDF pages~5--6, proves
\begin{equation}\label{eq:answer}
 \left\|\widehat A\right\|_{\frac{2(K-1)d}{(K-1)d+1}}
 \le C(d,K)\,\|A\|_{\op},
 \qquad
 C(d,K)\le K^{2d}\,
 \mathrm{BH}^{\le (K-1)d}_{\Omega_K}.
\end{equation}
Moreover, \(K^{2d}\) may be replaced by \(|\Sigma_K|^d\), where
\[
 \Sigma_K=\{(\ell,m)\in\Z_K^2:\ell\text{ and }m\text{ are coprime}\}.
\]
The cyclic-group BH constant on the right is finite independently of \(n\).
Thus \eqref{eq:answer} is genuinely dimension-free in the tensor length.

\section*{Idea of the proof in the answering paper}
For prime \(K\), the nonzero cyclic subgroups of \(\Z_K^2\) form a clean
partition.  For composite \(K\), primitive-generated subgroups overlap, and an
HW element can have spectrum either \(\Omega_K\) or the shifted root set
\(\Omega_{2K}\setminus\Omega_K\).  The later paper accommodates both effects
by averaging suitable eigenprojections into density matrices \(\rho(\omega)\).

The scalar function
\[
 f_A(\omega)=\operatorname{tr}\!\bigl(A\rho(\omega)\bigr)
\]
has degree at most \((K-1)d\).  Distinct HW profiles contribute disjoint sets
of scalar monomials, so their coefficients cannot cancel across profiles; the
averaging costs at most a factor \(|\Sigma_K|^d\).  The cyclic-group BH
inequality at degree \((K-1)d\), together with
\(|\operatorname{tr}(A\rho)|\le\|A\|_{\op}\), yields \eqref{eq:answer}.

\section*{Exact scope}
The source asked broadly what can be said in composite local dimension, and
\eqref{eq:answer} supplies a uniform theorem for every non-prime \(K\ge4\).
So that residual is literature-already-answered.  The modification is
substantive: the prime-case scalar degree \(d\) becomes \((K-1)d\), and the
coefficient exponent changes accordingly from \(2d/(d+1)\) to
\(2(K-1)d/((K-1)d+1)\).  This note therefore makes no claim that the
prime-strength composite inequality holds.  Nor does \cite{answer} settle the
source's separate optimal-growth question about whether exponential dependence
on \(d\) is necessary.

\paragraph{Provenance.}
This is a literature-status identification, not new mathematics from the run.
Exact arXiv-id and core-keyword searches through 11 August 2026 found no prior
packet recording this precise source-to-answer match.

\begin{thebibliography}{9}
\raggedright
\bibitem{source}
O. Klein, J. Slote, A. Volberg, and H. Zhang,
\emph{Quantum and classical low-degree learning via a dimension-free Remez
inequality}, arXiv:2301.01438v3 (2023).  See the conclusion on PDF page~32.

\bibitem{answer}
J. Slote, A. Volberg, and H. Zhang,
\emph{Noncommutative Bohnenblust--Hille Inequality for qudit systems},
arXiv:2406.08509 (2024).  See Theorem~7 on PDF pages~5--6 and its proof in
Section~3.
\end{thebibliography}

\end{document}
